% !TEX program = xelatex
\documentclass[11pt,a4paper]{article}

\usepackage{fontspec}
\usepackage{polyglossia}
\usepackage{amsmath,amssymb,amsfonts}
\usepackage{xcolor}
\usepackage{graphicx}
\usepackage{geometry}

\geometry{left=2cm,right=2cm,top=2cm,bottom=2cm}

\setdefaultlanguage{persian}
\setotherlanguage{english}
\setmainfont{Vazir}

\newcommand{\en}[1]{\textenglish{#1}}
\newcommand{\separator}{\vspace{1em}\noindent\rule{\linewidth}{0.4pt}\vspace{1em}}
\newcommand{\thickseparator}{\vspace{1em}\noindent\rule{\linewidth}{1pt}\vspace{1em}}

\title{\textbf{پاسخ سوال ۱} \\[0.4cm] \small بررسی تغییرات \en{discount numbers ($d$)} در \en{Laplace smoothing} با افزایش اندازه پیکره و تعداد \en{bigrams}}
\author{محمد باغبان زاده \\ شماره دانشجویی: شناسهٔ دانشجو}
\date{}
\linespread{1.35}

\begin{document}

\maketitle
\thickseparator

\section{مفاهیم پایه: مدل‌های زبانی N-gram و مشکل صفر شدن احتمالات}
مدل‌های زبانی \textbf{N-gram} یکی از ساده‌ترین و پرکاربردترین روش‌ها برای تخمین احتمال دنباله‌ای از کلمات هستند. در یک مدل بایگرام (bigram)، احتمال یک کلمه با شرط روی کلمهٔ قبلی به صورت زیر محاسبه می‌شود:
\[
P(w_i | w_{i-1}) = \frac{C(w_{i-1} w_i)}{C(w_{i-1})}
\]
که در آن \(C(w_{i-1} w_i)\) شمارش (count) رخداد بایگرام \(w_{i-1} w_i\) در پیکرهٔ آموزشی و \(C(w_{i-1})\) شمارش unigram کلمهٔ \(w_{i-1}\) است.

مشکل عمدهٔ این تخمین‌گر بیشینه‌درست‌نمایی (MLE) آن است که اگر یک بایگرام در پیکرهٔ آموزشی دیده نشده باشد، احتمال آن صفر محاسبه می‌شود. در مرحلهٔ آزمون نیز ممکن است همان بایگرام ظاهر شود و کل احتمال جمله را صفر کند. برای رفع این مشکل از تکنیک‌های \textbf{هموارسازی (smoothing)} یا \textbf{تخفیف (discounting)} استفاده می‌شود.

\separator
\section{هموارسازی لاپلاس \en{(Laplace smoothing)} و مفهوم تخفیف (discount)}
ساده‌ترین روش هموارسازی، \textbf{لاپلاس \en{(یا add-one)}} است. در این روش به شمارش تمام \en{N-gram}ها یک واحد اضافه می‌شود. برای بایگرام‌ها:
\[
P_{\text{Laplace}}(w_n | w_{n-1}) = \frac{C(w_{n-1}w_n) + 1}{C(w_{n-1}) + V}
\]
که \(V\) اندازهٔ دایرهٔ واژگان (تعداد کلمات یکتای پیکره) است. مخرج به این دلیل \(V\) واحد افزایش می‌یابد که به هر unigram نیز یک شمارش اضافه شده است.

برای درک مفهوم \textbf{تخفیف \en{(discount number)}}، معمولاً از \textbf{شمارش تعدیل‌شده \en{(adjusted count)}} استفاده می‌کنیم. شمارش تعدیل‌شده برای یک بایگرام با شمارش اصلی \(c\) به گونه‌ای تعریف می‌شود که پس از تقسیم بر \(C(w_{n-1})\)، همان احتمال لاپلاس حاصل گردد:
\[
c^*(w_{n-1}w_n) = [C(w_{n-1}w_n) + 1] \times \frac{C(w_{n-1})}{C(w_{n-1}) + V}
\]
سپس \textbf{عدد تخفیف \en{(discount number)}} \(d_c\) به صورت نسبت شمارش تعدیل‌شده به شمارش اصلی تعریف می‌شود:
\[
d_c = \frac{c^*}{c} = \frac{C(w_{n-1}w_n) + 1}{C(w_{n-1}w_n)} \times \frac{C(w_{n-1})}{C(w_{n-1}) + V}
\]
به عبارت ساده‌تر:
\[
d_c = \frac{c+1}{c} \cdot \frac{C(w_{n-1})}{C(w_{n-1}) + V}
\]
این عدد نشان می‌دهد که چه نسبتی از شمارش اصلی پس از هموارسازی باقی می‌ماند. هر چه \(d_c\) به ۱ نزدیک‌تر باشد، تخفیف کمتری اعمال شده است.

\separator
\section{تأثیر افزایش اندازهٔ پیکره و افزایش تعداد بایگرام‌ها بر \(d_c\)}
حال با توجه به خواسته مساله:  
«با در نظر گرفتن پیکره \(A\) و دایره واژگان \(V\)، هرچه اندازهٔ پیکره را افزایش دهیم و ترکیب‌های بیشتری مثلاً برای \en{bigrams} ایجاد کنیم، چه تغییری در \en{discount numbers (\(d\))} هنگام \en{Laplace smoothing} رخ می‌دهد؟»

برای پاسخ، تغییرات هر یک از مؤلفه‌های فرمول \(d_c\) را در شرایط زیر تحلیل می‌کنیم:
\begin{itemize}
\item اندازهٔ پیکره (تعداد کل توکن‌ها) افزایش می‌یابد.
\item در نتیجه، شمارش هر بایگرام خاص (\(c = C(w_{n-1}w_n)\)) معمولاً افزایش می‌یابد (مگر اینکه بایگرام بسیار نادر باشد که در پیکره‌های بزرگتر نیز احتمالاً حداقل یک بار دیده می‌شود).
\item شمارش هر \en{unigram} (\(C(w_{n-1})\)) نیز به طور چشمگیری افزایش می‌یابد.
\item دایرهٔ واژگان \(V\) ممکن است کمی افزایش یابد، اما در پیکره‌های بزرگ رشد آن بسیار کندتر از رشد شمارش‌هاست و برای تحلیل می‌توان آن را تقریباً ثابت در نظر گرفت.
\end{itemize}

فرمول \(d_c\) را به صورت حاصلضرب دو ضریب می‌نویسیم:
\[
d_c = \underbrace{\frac{c+1}{c}}_{\text{ضریب اول}} \times \underbrace{\frac{C(w_{n-1})}{C(w_{n-1}) + V}}_{\text{ضریب دوم}}
\]

\subsection*{رفتار ضریب اول \(\frac{c+1}{c}\)}
با افزایش \(c\) (که با بزرگتر شدن پیکره رخ می‌دهد)، مقدار \(\frac{c+1}{c} = 1 + \frac{1}{c}\) به سمت ۱ کاهش می‌یابد. برای مثال:
\begin{itemize}
\item اگر \(c=1\)، ضریب برابر \(2\) است.
\item اگر \(c=10\)، ضریب برابر \(1.1\) است.
\item اگر \(c=100\)، ضریب برابر \(1.01\) است.
\end{itemize}
بنابراین این ضریب از اعداد بزرگ (در پیکرهٔ کوچک) به سمت ۱ میل می‌کند.

\subsection*{رفتار ضریب دوم \(\frac{C(w_{n-1})}{C(w_{n-1}) + V}\)}
با افزایش \(C(w_{n-1})\) (که با بزرگ شدن پیکره افزایش می‌یابد)، این کسر نیز به سمت ۱ میل می‌کند، زیرا:
\[
\frac{C}{C+V} = 1 - \frac{V}{C+V}
\]
هنگامی که \(C\) بسیار بزرگتر از \(V\) می‌شود، جملهٔ \(\frac{V}{C+V}\) به صفر نزدیک می‌شود و کسر به ۱ نزدیک می‌گردد.

\subsection*{نتیجه برای \(d_c\)}
هر دو ضریب با افزایش اندازهٔ پیکره به سمت ۱ حرکت می‌کنند. در نتیجه حاصلضرب آنها یعنی \(d_c\) نیز به سمت ۱ افزایش می‌یابد. به بیان دقیق‌تر:
\[
\lim_{C(w_{n-1}) \to \infty} \lim_{c \to \infty} d_c = 1
\]
این بدان معناست که \textbf{عدد تخفیف \(d\) افزایش می‌یابد و به ۱ نزدیک می‌شود}. از دیدگاه عملی، هرچه پیکره بزرگتر باشد، نیاز به تخفیف (کم کردن از شمارش‌های دیده شده) کمتر می‌شود، زیرا بیشتر بایگرام‌های محتمل در پیکره دیده شده‌اند و جرم احتمال کمتری باید به بایگرام‌های دیده‌نشده اختصاص یابد.

\separator
\section{مثال عددی برای روشن‌تر شدن مطلب}
فرض کنید پیکرهٔ کوچکی داریم با \(C(w_{n-1})=10\) و \(V=1000\) و یک بایگرام خاص با \(c=2\). در این صورت:
\[
d_c = \frac{2+1}{2} \times \frac{10}{10+1000} = 1.5 \times 0.0099 = 0.01485 \quad (\text{بسیار کوچک})
\]
حال پیکره را بسیار بزرگ می‌کنیم. فرض کنید \(C(w_{n-1})=10000\)، \(V=2000\) (کمی افزایش یافته) و همان بایگرام اکنون \(c=200\) بار دیده شده است:
\[
d_c = \frac{200+1}{200} \times \frac{10000}{10000+2000} = 1.005 \times 0.8333 = 0.8375
\]
می‌بینیم که \(d_c\) از \en{۰.۰۱۴} به \en{۰.۸۳۷} افزایش یافته است. با بزرگتر شدن بیشتر پیکره، \(d_c\) به ۱ نزدیک‌تر خواهد شد.

\separator
\section{پاسخ نهایی}
با توجه به تحلیل فرمول تخفیف در هموارسازی لاپلاس:
\[
d_c = \frac{c+1}{c} \times \frac{C(w_{n-1})}{C(w_{n-1}) + V}
\]
و با استدلال ریاضی ارائه شده:
\begin{itemize}
\item افزایش اندازهٔ پیکره باعث افزایش شمارش‌ها (\(c\) و \(C(w_{n-1})\)) می‌شود.
\item هر دو ضریب کسری به سمت ۱ میل می‌کنند.
\item در نتیجه، \textbf{عدد تخفیف \(d\) افزایش یافته و به ۱ نزدیک می‌شود}.
\end{itemize}
به عبارت دیگر، هرچه پیکره بزرگتر شود و ترکیب‌های بیشتری (بایگرام‌های متنوع‌تر) در آن ظاهر شوند، تأثیر هموارسازی لاپلاس کاهش می‌یابد و شمارش تعدیل‌شده \(c^*\) به شمارش اصلی \(c\) نزدیکتر می‌شود. بنابراین \textbf{تغییر در \en{discount numbers (\(d\))} به صورت افزایش (به سمت ۱) خواهد بود}.

\thickseparator

\end{document}